\chapter{Remaining FYP Plan}
\section{Remaining Scope}
Complete the broader stock research workflow, evaluate its contributions and document its operational boundaries. Brokerage execution remains excluded.
\section{Remaining Modules and Features}
Add audited company coverage, paired news/social features and divergence, a service layer, prompted explanations with constrained claims and optional verified student adaptation. Review the profile wording and the actual-close timing policy.
\section{Remaining Iterations}
Iteration 2 completes engine/service and the preregistered divergence comparison. Iteration 3 (Mid-2) adds the prompted explainer, free-prose claim binding, NFR/parity/degraded-operation measurement and acceptance preparation. Iteration 4 compares original and QLoRA-adapted Phi-3-mini with the template and performs broader evaluation and handover.
\section{Implementation Plan}
First audit additional source coverage and available-label boundaries. Freeze the divergence protocol before evaluating variants. Load a verified original student checkpoint, score its raw prompted baseline, expand filtered examples and preserve the held-out set before adaptation. Integrate the tested service and interface incrementally.
\section{Testing and Validation Plan}
Repeat the timing/canary controls for each scientific identity. Compare macro-F1, baseline lift, Brier/reliability and class distributions under shared splits. Measure raw and gated explanation quality. Inject missing feeds/thin history and pin a historical clock for serving parity. The automated workflow is \textbf{2 jobs $\times$ 2 OS = 4 checks, then 1 gated delivery job}; local counts remain distinct from classifier or participant evidence.
\section{Deployment Plan}
Prepare a container and an explicit target VM, measure the required environments and deploy the research service. Test installation from the documented package on a separate machine. Current CI delivers downloadable source/PDF artifacts; it does not deploy a production service.
\section{Tentative Timeline}
Use the proposal's relative checkpoints in Table~\ref{tab:engineering-dates}. Formal assessment appointments have not been supplied. Engine/divergence/service work precedes the Mid-2 prompted explainer, and QLoRA follows a working prompted baseline.
\section{Expected Final Deliverables}
Deliver integrated source/service, source-bound experiments and ablations, timing/calibration and explanation evaluations, measured performance/degraded-operation evidence, external installation instructions, participant records where completed, and final report/presentation. Retain a defensible null-result report if the learned engine has no useful advantage.
